\section{Analysis and Discussion}
\label{sec:analysis}

\paragraph{What the controls isolate.}
The number that isolates our contribution is not \crossdoc against \doccausal but
\crossdoc against \concatlink at identical connectivity. The concat literature
\citep{yasunaga2022linkbert,shi2024incontext,staniszewski2025structured,yasunaga2022dragon}
is consistent that co-locating related documents under ordinary attention buys most of
the \emph{general} long-context gain; that gain is available to \concatlink and
\concat too, and we do not claim it. What \concatlink cannot do is tell the model
\emph{where} the entitlement to read the target begins. On RepoBench perplexity the
link-gated mask beats its whole-document superset, and the fully permissive
\concat is the worst condition of the four (\Cref{sec:results:controls}), the
ordering K-BERT's visible-matrix ablation \citep{liu2020kbert} would predict: the
typed edge is doing work that more attention does not. We treat the controls as
first-class results, not a survival test.

\paragraph{Why the inference-time density result is about capability, not the thesis.}
The density experiments say that on code the cross-document benefit is mostly a
property of what the mask exposes at evaluation: a \doccausal-trained model given
\crossdoc grants at test time recovers nearly the whole effect, and training-time
density saturates early. Read as a statement about \emph{intelligence}---how much
better a model predicts when it can see the file it imports---this is a sharp finding:
the capability to exploit an adjacent, causally readable document is largely already
present in a document-isolated model, and what training with grants adds is modest.
It is not, however, a statement about what \crossdoc is \emph{for}. The purpose of
training with the link-gated mask is that the model learns, under maintained
causality and through the plain next-token objective, that emitting a reference is
what opens access to the referenced document---the behaviour the generation loop
exploits. A \doccausal-trained model can be \emph{handed} a target and read it; it has
no reason to have learned to \emph{ask} for one. Whether the trained model in fact
generates resolvable references at a higher rate, and uses what comes back, is the
measurement the design was built for and the one this paper does not yet contain: the
models here are too small to generate coherently enough for the test to be meaningful.
We flag it as the central open measurement rather than let the density result stand in
for it.

\paragraph{Why code margins are smaller than Wikipedia's.}
The Wikipedia HotpotQA $\Delta$ is an order of magnitude larger than the RepoBench
$\Delta$s. A bridge question requires a fact that lives only in the second article; a
cross-file completion is usually locally predictable from the import line and the
surrounding code, and the target file resolves the residual uncertainty about a name
or a signature. Community-pack perplexity, which averages over every token of every
document, is accordingly near-noise for code and is why the density experiments report
it only as a within-model, same-token difference; the use-line re-anchoring of the
internal ports exists to concentrate the score where the dependency is actually
consumed.

\paragraph{The Wikipedia community-pack measurement.}
On the community-pack instrument the Wikipedia graph does not follow the code line,
and an earlier reading of it produced a small negative that appears to depend on how
well the model already fits the text. We do not interpret that measurement here: the
HotpotQA contrast on the same checkpoints is large and positive, the two instruments
differ in what they score (a link-dependent answer versus every token of a linked
article), and the community-pack result on this corpus is under audit. Until it is
resolved we make no claim, in either direction, about generic held-out Wikipedia text.
% User (2026-08-23) suspects the wiki community_pack negative is a harness/corruption
% artifact; see TODOS.md. Do not promote either reading until re-measured.

\paragraph{RoPE without a per-document reset.}
We use global rotary positions over the packed sequence, so a granted target sits at a
packing-order-dependent offset from its linker. \citet{zhao2024analysing} reset
positions per document, but only because their intra-document mask forbids all
cross-document attention; once a grant crosses a boundary a clean reset is ill-defined,
since one key has one absolute position but is read by linkers at many distances. The
targets-first packing order keeps those distances short and correlated, which also
means the prediction that the advantage \emph{grows} with packing distance is not
cleanly testable in our setup. We therefore treat position handling as an empirical
question (does link utilization decay with distance?) rather than ablate a reset; that
diagnostic and a per-grant additive embedding are future work.

\paragraph{Training with the same mask versus an auxiliary structure loss.}
A reviewer may ask whether a link- or target-prediction loss would beat the plain
next-token objective, since structure masks without a structure objective leave gains
on the table in GraphCodeBERT \citep{guo2021graphcodebert} and EMAT
\citep{wu2022emat}. Our position is that \emph{trained} integration is what those
results show matters, and the train--inference mirror provides it: the identical mask,
detector, and matching key are learned through the existing loss and run unchanged at
generation. We add no auxiliary loss.

\paragraph{Limitations.}
\emph{Scale and fit.} All models are $\approx350$M parameters trained for a few billion
tokens; the literature predicts retrieval-style gains shrink with scale
\citep{khandelwal2020knnlm,borgeaud2022retro,taylor2022galactica}, and whether
edge-gated attention follows that trend is an open question our scaling runs are
designed to test, not one we preconclude. \emph{arXiv.} Papers are long relative to a
$32$k window, so citation targets are rarely co-packed and the arXiv models are
uninformative about the method; the modality needs sequence-parallel training.
\emph{Placebo on the headline arms.} The cross arm of HotpotQA and RepoBench sees
strictly more tokens than the flat arm; the internal ports carry a content-swapped
placebo that separates the right document from any document, and the headline arms do
not yet. \emph{Leakage.} HotpotQA is built on 2017 English Wikipedia and The Stack
overlaps public code; the paired same-token contrast cancels memorization of the
completion but not re-exposure of a memorized target, and a leakage-stratified
analysis is pending. \emph{Fired-subset reporting.} $\Delta$ is over items where the
detector fires, which correlates with clean titles and parseable imports.
\emph{Inference cost.} The generation loop recomputes the full forward at every token
with no key/value cache. \emph{Diversity.} The merged-corpus experiment was invalidated
by a dataloader bug and is being re-run; no diversity claim is made here.
\emph{Packing details.} Realized link density is below raw graph density (targets not
co-packed, the grant cap, cycles dropped by the DAG gate), which we quantify only
through the realized-grants axis; and grants beyond the cap are dropped by position,
not importance.
